

\section{L'Encodage Profond : Au-delà de la liste bilingue}
space{0.3cm}
oindent

La didactique du grec ancien traverse, depuis plusieurs décennies, une période de remise en question
fondamentale, souvent qualifiée de crise de légitimité pédagogique. Au cœur de cette turbulence se
trouve une interrogation sur l'efficacité des méthodes traditionnelles, héritées du XIXe siècle, qui
ont longtemps privilégié une approche philologique et analytique au détriment de l'acquisition d'une
compétence de lecture fluide. L'objet de cette section de thèse, intitulée \og L'Encodage Profond :
dépasser la liste bilingue pour lier le mot à une image ou un contexte\fg{}, est d'explorer les
fondements cognitifs et psycholinguistiques d'un changement de paradigme nécessaire : le passage
d'un apprentissage par traduction (codage-recodage) à un apprentissage par sémantisation directe.

Traditionnellement, l'acquisition du vocabulaire grec s'opère par la mémorisation de listes
bilingues, où chaque item lexical de la langue cible (L2) est apparié à un item de la langue
maternelle (L1). Cette approche, bien qu'économique en apparence, repose sur le postulat implicite
que la maîtrise lexicale équivaut à la capacité de fournir une traduction rapide. Or, les recherches
contemporaines en psychologie cognitive, notamment celles portant sur le lexique bilingue et la
mémoire, suggèrent que ce mode d'apprentissage induit un traitement superficiel de l'information,
limitant drastiquement la rétention à long terme et l'automatisation des processus de lecture.

Ce rapport se propose d'analyser en profondeur les mécanismes de l'\textbf{encodage profond},
concept central pour comprendre comment un apprenant peut construire un réseau lexical robuste en
grec ancien sans la médiation systématique du français. Nous examinerons comment les théories de la
cognition incarnée (\textit{Embodied Cognition}), du double codage (\textit{Dual Coding Theory}) et
de la charge d'implication (\textit{Involvement Load Hypothesis}) convergent pour valider des
pratiques pédagogiques alternatives telles que la méthode Polis, l'approche \textit{Per Se
Illustrata} ou le TPR (\textit{Total Physical Response}). Il s'agira de démontrer que la
sémantisation directe — c'est-à-dire l'établissement d'un lien neuronal direct entre le mot grec et
le concept ou l'image mentale qu'il désigne — n'est pas seulement une technique d'animation de
classe, mais une nécessité neurocognitive pour passer d'une connaissance déclarative instable à une
compétence procédurale durable.

L'analyse s'articulera autour de quatre axes majeurs : la critique cognitive du modèle associatif
traditionnel, les fondements théoriques de l'encodage profond, l'étude détaillée des dispositifs de
sémantisation directe (image, action, contexte), et enfin les implications neurodidactiques pour
l'enseignement du grec ancien comme langue \og vivante\fg{} de culture.



\subsection{Critique cognitive de la liste bilingue : le piège de l'association de mots}

Pour légitimer l'adoption de stratégies d'encodage profond, il est impératif de déconstruire les
mécanismes cognitifs sous-jacents à l'apprentissage par listes bilingues, qui constitue encore la
norme dans la majorité des manuels universitaires et scolaires. Cette approche, souvent perçue comme
\og naturelle\fg{} ou \og incontournable\fg{} pour une langue morte, s'avère être un frein
structurel à l'acquisition de la fluidité.



\subsection{L'hégémonie du modèle d'association de mots (word association model)}

La recherche sur le bilinguisme a longtemps cherché à modéliser l'organisation du lexique mental
chez les apprenants de langues secondes. Deux modèles concurrents, proposés initialement par Potter,
So, Von Eckardt et Feldman (1984), permettent de situer l'enjeu : le \textit{Word Association Model}
(WAM) et le \textit{Concept Mediation Model} (CMM).\cite{III-1-3-ref1}

Dans le Modèle d'Association de Mots (WAM), le lexique de la L2 est subordonné à celui de la L1.
Concrètement, pour un étudiant apprenant le mot grec ὁ ἵππος via une liste \og hippos \=
cheval\fg{}, le chemin neuronal pour accéder au sens est indirect. La perception du mot grec active
d'abord la représentation lexicale du mot français \og cheval\fg{}, qui elle-même active le concept
sémantique de l'animal équidé.

Le schéma d'accès est donc : L2 (Forme) $\\rightarrow$ L1 (Forme) $\\rightarrow$ Concept.

À l'inverse, le Modèle de Médiation Conceptuelle (CMM) postule un accès direct du mot L2 au concept,
sans passer par la L1.

Le schéma est alors : L2 (Forme) $\\rightarrow$ Concept $\\leftarrow$ L1 (Forme).

La liste bilingue, par sa structure même, force le cerveau de l'apprenant à consolider la route du
WAM. Chaque répétition de la paire \og grec-français\fg{} renforce le lien entre les deux étiquettes
verbales, mais ne construit pas de lien fort entre le mot grec et la réalité qu'il désigne. Kroll et
Stewart (1994), dans leur \textbf{Modèle Hiérarchique Révisé (Revised Hierarchical Model \- RHM)},
ont démontré que si les débutants s'appuient naturellement sur l'association lexicale, la véritable
compétence bilingue (fluency) se caractérise par une transition vers la médiation
conceptuelle.\cite{III-1-3-ref4} Le drame didactique du grec ancien réside dans le fait que les
méthodes traditionnelles, en maintenant l'usage intensif de la traduction et des listes jusqu'à des
niveaux avancés, \og fossilisent\fg{} les apprenants dans ce stade débutant. L'étudiant, même après
plusieurs années, continue de traduire mentalement chaque mot, ce qui sature sa mémoire de travail
et empêche une lecture fluide et intuitive.



\subsection{La pauvreté du traitement : l'illusion de l'apprentissage superficiel}

Au-delà de l'architecture du lexique, la liste bilingue pose un problème de \og profondeur\fg{} de
traitement. La théorie des \textbf{Niveaux de Traitement (Levels of Processing)}, formulée par Craik
et Lockhart (1972), conteste l'idée que la simple répétition suffit à la mémorisation à long terme.
Selon cette théorie, la durabilité de la trace mnésique dépend de la profondeur de l'analyse
effectuée lors de l'encodage.\cite{III-1-3-ref7}

\begin{itemize}\item \textbf{Traitement Superficiel (Shallow Processing) :} Il se focalise sur les caractéristiques physiques ou sensorielles du stimulus. Apprendre une liste de vocabulaire implique souvent un traitement visuel (orthographe) ou phonologique (sonorité) de la paire de mots, sans nécessairement engager une réflexion sémantique intense. La répétition mécanique (\og par cœur\fg{}) est l'archétype de ce traitement de surface, souvent appelé \og répétition de maintenance\fg{} (\textit{maintenance rehearsal}). Elle maintient l'information active en mémoire à court terme mais échoue à l'intégrer durablement en mémoire à long terme.\cite{III-1-3-ref8}\item \textbf{Traitement Profond (Deep Processing) :} Il implique une analyse sémantique, l'élaboration de liens avec des connaissances antérieures, ou la création d'images mentales. C'est ce type de traitement qui permet de consolider l'information.\cite{III-1-3-ref11}\end{itemize}En didactique du grec, l'usage exclusif de la liste bilingue favorise un traitement superficiel.
L'étudiant peut \og savoir\fg{} que βαίνω signifie \og je marche\fg{} pour l'interrogation du
lendemain, mais cette connaissance est fragile car elle n'est connectée à aucun réseau sémantique
riche (mouvement, espace, corps) autre que l'étiquette française. Le mot reste une coquille vide,
une simple correspondance formelle.



\subsection{La surcharge cognitive et l'effet d'attention partagée}

Un troisième facteur critique est la gestion de la charge cognitive. La \textbf{Théorie de la Charge
Cognitive (Cognitive Load Theory)} de Sweller met en évidence les limites de la mémoire de travail.
L'apprentissage par listes ou par gloses bilingues déportées (en fin de volume ou en bas de page)
induit un \textbf{Effet d'Attention Partagée (Split Attention Effect)}.\cite{III-1-3-ref13}

Lorsque l'apprenant lit un texte grec et doit constamment détourner son regard pour consulter une
liste de vocabulaire, il doit intégrer mentalement deux sources d'information disparates (le texte
et la définition) séparées spatialement et temporellement. Cet effort d'intégration consomme des
ressources cognitives précieuses, créant une charge cognitive \og extrinsèque\fg{} qui n'est pas
liée à la complexité intrinsèque du grec, mais à la mauvaise conception du matériel
pédagogique.\cite{III-1-3-ref15}

Au lieu de consacrer son énergie mentale à construire le sens du texte (la sémantisation),
l'étudiant la gaspille dans des opérations de recherche visuelle et de maintien en mémoire de
travail des correspondances lexicales. Cette fragmentation de l'attention empêche l'immersion et la
construction d'un modèle mental cohérent de la situation décrite par le texte.

Le tableau suivant synthétise les déficits cognitifs inhérents à l'approche par liste bilingue :

\begin{table}[h!]\small\centering\begin{tabular}{| p{2.3cm} | p{3.5cm} | p{3.5cm} |}\hline\textbf{Dimension Cognitive} & \textbf{Caractéristique de la Liste Bilingue} & \textbf{Conséquence pour l'Apprenant de Grec Ancien} \\ \hline\textbf{Accès Lexical} & Indirect (médiation par la L1) & Lenteur de lecture, impossibilité de penser en langue cible. \\ \hline\textbf{Profondeur de Traitement} & Superficielle (répétition de maintenance) & Oubli rapide, connaissance déclarative non transférable. \\ \hline\textbf{Charge Cognitive} & Élevée (Split Attention) & Fatigue mentale, saturation de la mémoire de travail, décrochage. \\ \hline\textbf{Contextualisation} & Nulle (mots isolés) & Incapacité à saisir la polysémie ou les nuances sémantiques réelles. \\ \hline\end{tabular}\end{table}

\subsection{Fondements théoriques de l'encodage profond}

Face aux impasses du modèle associatif, l'\textbf{encodage profond} propose une alternative fondée
sur la multiplication des connexions neuronales lors de l'apprentissage. Lier le mot à une image,
une action ou un contexte ne relève pas de l'accessoire ludique, mais d'une stratégie cognitive
puissante visant à la \og sémantisation directe\fg{}.



\subsection{La théorie du double codage (dual coding theory) : la puissance de l'image}

La \textbf{Théorie du Double Codage (DCT)}, développée par Allan Paivio, est le socle théorique le
plus solide pour justifier l'usage de l'image en didactique des langues anciennes. Selon Paivio, le
système cognitif humain traite l'information via deux canaux distincts mais interconnectés :

1. \textbf{Le système verbal (logogens) :} Spécialisé dans le traitement du langage, des sons et des
structures séquentielles.

2. \textbf{Le système non-verbal ou imagé (imagens) :} Spécialisé dans le traitement des objets
visuels, des scènes et des relations spatiales.\cite{III-1-3-ref18}

L'apport crucial de la DCT réside dans le concept d'\textbf{additivité}. Lorsqu'un mot est appris en
association avec une image (par exemple, le mot grec ἡ ναῦς présenté avec l'image d'un navire
antique), l'information est encodée simultanément dans les deux systèmes. Cette redondance crée une
trace mnésique doublement robuste. Si le code verbal devient inaccessible (l'étudiant oublie le
mot), le code imagé peut servir de voie de récupération (\textit{retrieval cue}) pour réactiver le
mot, et inversement.\cite{III-1-3-ref20}

C'est ce que la recherche nomme l'\textbf{Effet de Supériorité de l'Image (Picture Superiority
Effect)}. Les stimuli picturaux sont généralement mieux rappelés que les stimuli verbaux car ils
sont plus susceptibles d'être encodés doublement.\cite{III-1-3-ref22} Pour le grec ancien, cela
signifie que l'illustration dans un manuel comme \textit{Alexandros} n'est pas une simple décoration
; elle est un vecteur d'encodage. Elle permet de court-circuiter la L1 : au lieu de lier ναῦς à \og
navire\fg{}, l'apprenant lie ναῦς directement à l'image mentale du bateau, initiant ainsi la
médiation conceptuelle directe décrite par Kroll.



\subsection{L'hypothèse de la charge d'implication (involvement load hypothesis)}

Si l'image fournit un support sensoriel, comment s'assurer que l'apprenant traite l'information en
profondeur? C'est ici qu'intervient l'\textbf{Hypothèse de la Charge d'Implication (Involvement Load
Hypothesis \- ILH)} proposée par Batia Laufer et Jan Hulstijn (2001). Cette théorie postule que
l'efficacité de l'acquisition du vocabulaire dépend du degré d'engagement cognitif et motivationnel
induit par la tâche.\cite{III-1-3-ref24}

La charge d'implication se compose de trois facteurs :

1. \textbf{Le Besoin (Need) :} La motivation extrinsèque ou intrinsèque à comprendre le mot. Le
besoin est \og modéré\fg{} si le professeur demande d'apprendre le mot, mais \og fort\fg{} si la
compréhension du mot est indispensable pour saisir le sens d'une histoire ou résoudre un
problème.\cite{III-1-3-ref27}

2. \textbf{La Recherche (Search) :} L'effort cognitif déployé pour trouver le sens du mot. Consulter
une glose en marge constitue une recherche nulle ou faible. En revanche, devoir inférer le sens d'un
mot à partir du contexte narratif ou de l'analyse d'une image constitue une recherche forte.

3. \textbf{L'Évaluation (Evaluation) :} L'effort fait pour comparer le sens trouvé avec le contexte
(ex: \og Est-ce que 'courir' a du sens ici, ou est-ce plutôt 's'enfuir'?\fg{}). C'est une prise de
décision sémantique.\cite{III-1-3-ref28}

L'application de l'ILH à la didactique du grec est radicale. Elle condamne les listes bilingues
(Need faible, Search nul, Evaluation nulle) et plébiscite les approches contextuelles inductives.
Dans une méthode comme \textit{Athenaze} (utilisée sans traduction) ou \textit{Polis}, l'étudiant
rencontre un mot inconnu et doit mobiliser des indices (image, racine étymologique, contexte de la
phrase précédente) pour en déduire le sens. Cet effort de \textbf{Recherche} et
d'\textbf{Évaluation} augmente considérablement la charge d'implication, garantissant un encodage
plus profond et une meilleure rétention. Comme le résume l'adage cognitif implicite de cette théorie
: \og Ce qui est facile à décoder est facile à oublier ; ce qui demande un effort de sémantisation
est durable\fg{}.



\subsection{La cognition incarnée (embodied cognition) et l'ancrage sensorimoteur}

Au-delà du visuel et du contextuel, l'encodage profond peut être \textbf{kinesthésique}. Les
théories de la \textbf{Cognition Incarnée (Embodied Cognition)} remettent en cause la vision
computationnelle de l'esprit (le cerveau comme ordinateur traitant des symboles abstraits) pour
affirmer que la cognition est fondamentalement enracinée dans les interactions sensorimotrices du
corps avec l'environnement.\cite{III-1-3-ref29}

Selon cette perspective, comprendre le verbe \og lancer\fg{} (βάλλω) implique une simulation
mentale, inconsciente et ultra-rapide, des circuits moteurs impliqués dans l'action de lancer. Le
sens n'est pas une définition de dictionnaire stockée dans une base de données abstraite, mais une
réactivation partielle de l'expérience corporelle.

Cette théorie offre une base scientifique solide à la méthode TPR (Total Physical Response) de James
Asher. En demandant aux apprenants d'exécuter physiquement des ordres en grec (ἀνάστηθι \- lève-toi,
περιπάτει \- marche), on exploite la synergie entre le système moteur et le système langagier.
L'encodage devient multimodal : le mot est lié au son (auditif), au concept (sémantique) et au
mouvement (proprioceptif). Cette \og trace mnésique incarnée\fg{} est particulièrement résistante à
l'oubli et favorise une récupération rapide, essentielle pour la fluidité orale et la
lecture.\cite{III-1-3-ref31}



\subsection{Le modèle déclaratif/procédural de ullman : vers l'automatisation}

Enfin, pour comprendre l'objectif ultime de la sémantisation directe, il faut se référer au
\textbf{Modèle Déclaratif/Procédural} de Michael Ullman. Ce modèle neurocognitif distingue deux
systèmes de mémoire pour le langage :

1. \textbf{La Mémoire Déclarative :} Gérée par le lobe temporal médian (hippocampe), elle stocke les
faits et événements explicites, ainsi que le lexique arbitraire. C'est une mémoire \og savoir
que\fg{}, consciente et rapide à acquérir mais demandant de l'attention pour être récupérée.

2. \textbf{La Mémoire Procédurale :} Gérée par les structures fronto-striatales (ganglions de la
base), elle gère les habiletés motrices et cognitives, ainsi que la grammaire et les séquences
automatisées. C'est une mémoire \og savoir comment\fg{}, implicite et
automatique.\cite{III-1-3-ref34}

L'apprentissage traditionnel par listes stocke le vocabulaire en mémoire déclarative. L'apprenant
doit \og réfléchir\fg{} pour retrouver le mot. L'objectif de l'encodage profond, via la répétition
en contexte (LSE, TPR, lecture extensive), est de transférer cette compétence vers la mémoire
procédurale ou, à tout le moins, d'automatiser l'accès lexical (\textit{lexical
access}).\cite{III-1-3-ref37} La sémantisation directe vise à ce que la perception du mot δένδρον
déclenche instantanément l'image de l'arbre, un processus réflexe proche du procédural, libérant
ainsi les ressources déclaratives pour la compréhension syntaxique complexe ou l'analyse littéraire.



\subsection{La sémantisation directe en pratique : dispositifs et méthodes}

Les fondements théoriques exposés ci-dessus ne sont pas de simples spéculations ; ils trouvent une
incarnation concrète dans plusieurs approches didactiques contemporaines qui rompent avec le modèle
grammatical. Cette section analyse comment ces méthodes opérationnalisent l'encodage profond pour
enseigner le grec ancien.



\subsection{Le paradigme "per se illustrata" : l'image comme pivot sémantique}

L'influence de Hans Ørberg et de sa méthode pour le latin (\textit{Lingua Latina Per Se Illustrata})
est déterminante. Bien que conçue pour le latin, son adaptation au grec, notamment à travers
l'ouvrage \textbf{\og Alexandros: To Hellenikon Paidion\fg{}} de Mario Díaz Ávila, illustre
parfaitement la sémantisation directe par l'image.\cite{III-1-3-ref39}



\subsection{Mécanisme de l'induction illustrée}

Dans \textit{Alexandros}, le texte est monolingue (entièrement en grec ancien). Le sens n'est jamais
donné par une traduction, mais construit par la triangulation : \textbf{Texte \- Image \- Contexte}.

\begin{itemize}\item \textbf{L'Iconicité Définitionnelle :} Contrairement aux manuels classiques où l'image est culturelle ou décorative (une photo de vase grec à côté d'un texte sur Socrate), les illustrations dans \textit{Alexandros} sont fonctionnelles. Si le texte dit ὁ ἥλιος λάμπει (le soleil brille), l'image montre un soleil brillant. L'image sert de définition instantanée, activant le système \textit{imagen} de Paivio.\cite{III-1-3-ref42}\item \textbf{Contiguïté Spatiale :} Les illustrations sont placées en marge immédiate du texte ou insérées dans le fil narratif. Cela respecte le principe de contiguïté de Mayer et évite l'effet d'attention partagée.\cite{III-1-3-ref14} L'œil capture le mot et l'image dans une même saccade visuelle, favorisant l'association neuronale simultanée.\item \textbf{L'Expansion Lexicale Endogène :} Une fois un noyau de vocabulaire concret acquis par l'image, les mots abstraits sont définis par des synonymes ou des antonymes grecs (ex: μικρός est défini comme le contraire de μέγας, illustré par un grand et un petit objet). L'apprenant reste dans le \og bain\fg{} linguistique grec, renforçant le réseau sémantique interne à la L2.\cite{III-1-3-ref44}\end{itemize}

\subsection{La méthode polis et le concept de "sémantisation discursive"}

Christophe Rico, linguiste et fondateur de l'Institut Polis à Jérusalem, a poussé la logique de
l'encodage profond plus loin en théorisant la \og sémantisation\fg{} comme un processus discursif et
immersif. Sa méthode, \textit{Polis: Speaking Ancient Greek as a Living Language}, repose sur le
refus radical de la traduction comme outil d'apprentissage initial.\cite{III-1-3-ref46}



\subsection{Définition de la sémantisation chez rico}

Pour Rico, la sémantisation n'est pas le simple étiquetage d'un objet. C'est l'acte par lequel
l'apprenant infère le sens probable d'un signe linguistique à partir de la situation d'énonciation.
Il distingue la conceptualisation linguistique (le sens du mot dans la langue) de la
conceptualisation discursive (le sens du mot en contexte).\cite{III-1-3-ref49}

En classe Polis, l'enseignant ne dit pas \og Ceci est une pierre\fg{}. Il prend une pierre, la
montre, la soupèse, la lance peut-être, en disant ὁ λίθος. L'apprenant doit connecter l'objet
physique, le geste de l'enseignant et le son entendu. Cette sémantisation directe empêche
l'interférence de la L1 et crée une représentation conceptuelle riche, incluant les propriétés
physiques de l'objet (poids, texture) qui sont absentes de la simple traduction écrite \og
pierre\fg{}.\cite{III-1-3-ref50}



\subsection{L'intégration du tpr (total physical response)}

Rico intègre systématiquement le TPR pour le vocabulaire d'action et spatial. Les étudiants doivent
réagir physiquement aux impératifs grecs.

\begin{itemize}\item \textbf{Exemple :} L'enseignant dit βάλε τὸν λίθον ἐπὶ τῆς κεφαλῆς (mets la pierre sur la tête). L'étudiant exécute.\item \textbf{Analyse :} Cette pratique valide la compréhension sans passer par la verbalisation en L1. Elle réduit l'anxiété (filtre affectif) et engage la mémoire procédurale motrice.\cite{III-1-3-ref50}\end{itemize}

\subsection{Le living sequential expression (lse) : l'encodage procédural}

Une innovation spécifique à la méthode Polis, et particulièrement pertinente pour la thèse, est le
\textbf{Living Sequential Expression (LSE)}. Inspirée des \og Séries\fg{} de François Gouin (XIXe
siècle), cette technique consiste à enseigner le vocabulaire non pas par mots isolés, mais par
séquences d'actions logiquement connectées.\cite{III-1-3-ref50}

\begin{itemize}\item \textbf{Le principe de la Série :} Au lieu d'apprendre \og feu\fg{}, \og allumette\fg{}, \og brûler\fg{} séparément, l'étudiant apprend une séquence : \og Je prends une allumette, je la gratte, la flamme jaillit, j'allume le bois, le feu brûle\fg{}.\item \textbf{Analyse Cognitive :} Le cerveau humain est particulièrement apte à mémoriser des scripts ou des schémas événementiels (\textit{event schemas}). Le LSE structure le lexique en chaînes causales et temporelles. Selon le modèle d'Ullman, cette structure séquentielle favorise l'intégration en mémoire procédurale (qui gère les séquences), contrairement aux listes qui saturent la mémoire déclarative.\cite{III-1-3-ref47} L'apprenant \og vit\fg{} la séquence (réellement ou par mime), ce qui constitue une forme avancée de sémantisation contextuelle.\end{itemize}

\subsection{Tprs et la narration comme vecteur sémantique}

Enfin, le \textbf{TPRS (Teaching Proficiency through Reading and Storytelling)}, bien que né pour
les langues vivantes, est adapté au grec ancien pour lier le mot à un contexte émotionnel et
narratif.

\begin{itemize}\item \textbf{Mécanisme :} Le vocabulaire est introduit au service d'une histoire co-construite avec les élèves. Les mots sont répétés de multiples fois (\textit{repetition flood}) mais dans des contextes variés et souvent humoristiques ou absurdes.\item \textbf{Apport :} La composante émotionnelle et narrative active la mémoire épisodique. L'apprenant se souvient du mot πίπτω (tomber) non pas par sa définition, mais parce qu'il se souvient du moment où le personnage de l'histoire est tombé de manière ridicule. Le \og Besoin\fg{} (Need) de comprendre l'histoire pour rire ou suivre l'intrigue maximise la charge d'implication.\cite{III-1-3-ref47}\end{itemize}Le tableau ci-dessous résume la comparaison entre l'approche traditionnelle et l'approche par
encodage profond :

\begin{table}[h!]\small\centering\begin{tabular}{| p{2cm} | p{2.2cm} | p{2.2cm} | p{2.2cm} |}\hline\textbf{Critère} & \textbf{Approche Liste Bilingue (Traditionnelle)} & \textbf{Approche Encodage Profond (Sémantisation Directe)} & \textbf{Fondement Théorique} \\ \hline\textbf{Accès au Sens} & Indirect (via L1) & Direct (L2 $\\rightarrow$ Concept) & RHM (Kroll \& Stewart) \\ \hline\textbf{Support Mnésique} & Paires verbales isolées & Image, Geste, Séquence narrative & Dual Coding (Paivio) \\ \hline\textbf{Rôle de l'Apprenant} & Passif (Mémorisation par cœur) & Actif (Infèrence, Action, Construction) & Involvement Load (Laufer) \\ \hline\textbf{Mémoire Ciblée} & Déclarative (Lexique explicite) & Procédurale / Épisodique (Automatismes) & Modèle Ullman \\ \hline\textbf{Contextualisation} & Faible ou nulle & Forte (Phrases, Récits, Séries) & Embodied Cognition \\ \hline\end{tabular}\end{table}

\subsection{Limites et perspectives : le défi de l'abstraction}

Si l'encodage profond par l'image et l'action est redoutablement efficace pour le vocabulaire
concret, une objection fréquente concerne le vocabulaire abstrait, omniprésent dans la littérature
grecque (philosophie, théologie). Comment \og sémantiser directement\fg{} des concepts comme ἀρετή
(vertu/excellence) ou λόγος (parole/raison)?



\subsection{L'abstraction par la contextualisation narrative}

La réponse des didacticiens de l'encodage profond réside dans l'extension du concept d'image.
L'image n'est pas seulement picturale, elle peut être situationnelle. Pour enseigner δικαιοσύνη
(justice), on ne montre pas une statue de la justice, on raconte une histoire (parabole, mythe) où
un personnage agit de manière juste et un autre de manière injuste. Le concept est encodé via une
\og scène sémantique\fg{} complexe.\cite{III-1-3-ref49}

C'est ici que la méthode inductive d'Alexandros ou les cours de Polis montrent leur pertinence : le
vocabulaire abstrait est introduit tardivement, une fois que la compétence de lecture permet de
comprendre des définitions ou des histoires complexes en grec. La sémantisation se fait alors par
synonymie, antonymie et paraphrase en L2, maintenant le principe de la médiation conceptuelle
directe.\cite{III-1-3-ref57}



\subsection{Le rôle métalinguistique de la l1}

Il ne s'agit pas de bannir totalement la L1, mais de changer son statut. Dans l'encodage profond, la
L1 ne sert pas à \textit{apprendre} le mot (encodage), mais éventuellement à \textit{vérifier} ou
\textit{affiner} la compréhension a posteriori (métacognition). L'objectif est que la représentation
mentale première du mot grec soit conceptuelle, la traduction française n'étant qu'une étiquette
secondaire disponible si nécessaire, mais non indispensable à la lecture
courante.\cite{III-1-3-ref48}



\subsection{Conclusion}

Nous montrons une rupture épistémologique et
pédagogique majeure. L'abandon de la liste bilingue au profit de stratégies d'encodage profond
(image, TPR, contexte narratif, LSE) n'est pas une concession à la facilité ou au ludique, mais une
exigence de rigueur cognitive.

Les théories psycholinguistiques du double codage, de la charge d'implication et de la cognition
incarnée convergent pour démontrer que le cerveau apprend mieux lorsqu'il multiplie les voies
d'accès au sens et qu'il est activement engagé dans la construction de la signification.

En didactique du grec ancien, cela implique une refonte des supports (vers des manuels type Per Se
Illustrata ou Alexandros) et des pratiques de classe (vers l'oralité et l'action type Polis). La
sémantisation directe est la clé pour transformer le grec ancien d'un code mort à déchiffrer en une
langue vivante de culture, accessible par une lecture fluide et intuitive. C'est en liant le mot à
l'image et au geste que l'on permet à l'étudiant de ne plus seulement traduire les Grecs, mais de
commencer, modestement, à penser avec eux.

Références Bibliographiques Intégrées (Sélection indicative) :

.\cite{III-1-3-ref1}

